\section{Penghapusan Left Recursion}
\label{sec:ch09-left-recursion}

\subsection{Mengapa Left Recursion Bermasalah}

Parser top-down seperti recursive descent memanggil fungsi untuk
nonterminal. Jika aturan pertama berbentuk \(A \to A\alpha\), fungsi
\code{parseA} akan memanggil dirinya sendiri sebelum mengonsumsi
token, sehingga terjadi rekursi tak berujung. Left recursion harus
dihapus sebelum parsing top-down.

\subsection{Left Recursion Langsung}

Aturan berbentuk
\[
A \to A\alpha_1 \mid A\alpha_2 \mid \beta_1 \mid \beta_2
\]
dapat diubah menjadi
\[
A \to \beta_1 A' \mid \beta_2 A', \qquad
A' \to \alpha_1 A' \mid \alpha_2 A' \mid \varepsilon.
\]
Nonterminal \(A'\) menangani pengulangan, sedangkan \(A\) memilih
alternatif dasar.

\begin{contoh}[Ekspresi aritmetika]
Aturan
\[
E \to E + T \mid T
\]
diubah menjadi
\[
E \to T E', \qquad E' \to + T E' \mid \varepsilon.
\]
Pohon urai tetap sama, tetapi parser dapat memutuskan berdasarkan
token pertama.
\end{contoh}

\begin{figure}[H]
\centering
\begin{tikzpicture}[node distance=6mm and 12mm]
  \node[phase, text width=4.4cm] (before)
        {Sebelum: \(E \to E + T \mid T\)};
  \node[phase, text width=4.4cm, right=of before] (after)
        {Sesudah: \(E \to T E'\)};
  \node[phase, text width=4.4cm, below=of after] (tail)
        {\(E' \to + T E' \mid \varepsilon\)};
  \draw[arrow] (before) -- (after);
  \draw[arrow] (after) -- (tail);
\end{tikzpicture}
\caption{Transformasi left recursion langsung}
\label{fig:ch09-left-recursion}
\end{figure}

\subsection{Left Recursion Tidak Langsung}

Left recursion tidak langsung terjadi ketika \(A\) dapat menurunkan
\(A\) melalui beberapa nonterminal:
\[
A \to B\alpha, \qquad B \to A\beta.
\]
Algoritma umum:
\begin{enumerate}[label=\arabic*.,leftmargin=*]
  \item urutkan nonterminal \(A_1, \ldots, A_n\);
  \item untuk setiap \(A_i\), ganti aturan yang dimulai dengan
        \(A_j\) (\(j < i\)) menggunakan aturan \(A_j\);
  \item hapus left recursion langsung pada \(A_i\).
\end{enumerate}

Proses ini dapat mengubah bentuk grammar secara signifikan. Setelah
transformasi, periksa kembali bahasa yang dibangkitkan dan uji dengan
beberapa string.

\subsection{Kondisi Nullable}

Jika \(\alpha\) dapat menurunkan \(\varepsilon\), transformasi harus
hati-hati. Aturan \(A \to A\alpha\) dengan \(\alpha\) nullable dapat
menyebabkan loop epsilon. Pastikan tidak ada siklus
\(\varepsilon\) dan periksa FIRST/FOLLOW setelah transformasi.

\begin{infobox}
Left recursion bukan kesalahan sintaksis. Banyak grammar alami
menggunakan left recursion karena mencerminkan associativity kiri.
Transformasi hanya diperlukan untuk parser top-down; parser
bottom-up seperti LR justru dapat menanganinya secara langsung.
\end{infobox}
